\documentclass[../../main.tex]{subfiles}

\begin{document}

\section{紧算子}\label{section:ch16-compact}

\begin{definition}[紧算子]\label{definition:ch16-compact}
算子 $T\in\mathcal{L}(H)$ 称为紧的, 若 $T(B)$ 关于强拓扑有紧的闭包, 其中 $B$ 是 $H$ 中的闭单位球.
\end{definition}

任何算子 $T\in\mathcal{L}(H)$ 将有界集映到有界集. 算子 $T\in\mathcal{L}(H)$ 是有限秩的, 若它的象是有限维的. 由于有限维空间的有界子集有紧闭包, 每个有限秩算子是紧的. 特别地, 若 $H$ 是有限维的, 则 $\mathcal{L}(H)$ 中的每个算子是紧的. 另一方面, 根据 Riesz 定理(定理 11), 若 $H$ 是无穷维的, 恒等算子 $\operatorname{Id}:H\to H$ 不是紧的. 同理, 若 $H$ 是无穷维的, $\mathcal{L}(H)$ 中的可逆算子不是紧的.

在任何度量空间, 集合的紧性与序列紧性是相同的. 此外, 由于度量空间是紧的当且仅当它是完备与全有界的, 完备度量空间的子集有紧闭包当且仅当它是全有界的. 我们因此有以下关于有界线性算子紧性的有用刻画.

\begin{proposition}\label{proposition:ch16-compact-equiv}
令 $H$ 为 Hilbert 空间而 $K$ 属于 $\mathcal{L}(H)$. 则以下断言等价:
\begin{enumerate}[(i)]
\item $K$ 是紧的.
\item $K(B)$ 是全有界的, 其中 $B$ 是 $H$ 中的闭单位球.
\item 若 $\{h_{n}\}$ 是 $H$ 中的有界序列, 则 $\{K(h_{n})\}$ 有强收敛的子序列.
\end{enumerate}
\end{proposition}

\begin{example}\label{example:ch16-diagonal-compact}
令 $\{\varphi_{k}\}$ 为可分 Hilbert 空间 $H$ 的规范正交基, 而 $\{\lambda_{k}\}$ 是收敛于 $0$ 的实数序列. 对 $h\in H$ 定义
\begin{align*}
T(h)=\sum_{k=1}^{\infty}\lambda_{k}\langle h,\varphi_{k}\rangle\varphi_{k}.
\end{align*}
我们从 Bessel 不等式与 $\{\lambda_{k}\}$ 的有界性推出 $T$ 属于 $\mathcal{L}(H)$ 且我们宣称 $T$ 是紧的. 根据前一个命题, 为证明 $T$ 是紧的仅须证明 $T(B)$ 是全有界的. 令 $\varepsilon>0$. 选取 $N$ 使得对 $k\geqslant N$, $|\lambda_{k}|<\varepsilon/2$. 定义 $T_{N}\in\mathcal{L}(H)$ 为
\begin{align*}
\text{对 }h\in H,\quad T(h)=\sum_{k=1}^{N}\lambda_{k}\langle h,\varphi_{k}\rangle\varphi_{k}.
\end{align*}
我们从 Bessel 不等式推出对 $h\in H$, $\|T(h)-T_{N}(h)\|<\varepsilon/2\|h\|$. 但 $T_{N}(B)$ 是有限维空间的有界子集, 因此它是全有界的. 令 $\varepsilon>0$. 存在 $T_{N}(B)$ 的有限 $\varepsilon/2$ 网而对该网的每个球的半径加倍我们得到 $T(B)$ 的有限 $\varepsilon$ 网.
\end{example}

线性算子 $T:H\to H$ 属于 $\mathcal{L}(H)$ 当且仅当它将弱收敛序列映到弱收敛序列(见习题\cref{exercise:ch16-47}).

\begin{proposition}\label{proposition:ch16-compact-weak-strong}
令 $H$ 为 Hilbert 空间. 则 $\mathcal{L}(H)$ 中的算子 $T$ 是紧的当且仅当它将弱收敛序列映到强收敛序列, 即
\begin{align*}
\text{若 }\{h_{n}\}\rightharpoonup h,\quad\text{则 }\{T(h_{n})\}\to T(h).
\end{align*}
\end{proposition}

\begin{proof}
根据前一个命题, 算子是紧的当且仅当它将有界序列映到有强收敛子序列的序列. 首先假设 $T$ 是紧的. 观察到对任何算子 $T\in\mathcal{L}(H)$, 若 $\{u_{k}\}\rightharpoonup u$, 则 $\{T(u_{k})\}\rightharpoonup T(u)$, 这是由于对每个 $v\in H$,
\begin{align*}
\lim_{k\to\infty}\langle T(u_{k}),v\rangle=\lim_{k\to\infty}\langle u_{k},T^{*}(v)\rangle=\langle u,T^{*}(v)\rangle=\langle T(u),v\rangle.
\end{align*}
令 $\{h_{n}\}\rightharpoonup h$. 根据 $T$ 的紧性, $\{T(h_{n})\}$ 的每个子序列有强收敛的子序列, 根据先前的观察, 它的强极限必须是 $T(h)$. 因此整个序列 $\{T(h_{n})\}$ 强收敛于 $T(h)$. 为证明反命题, 假设 $T$ 将弱收敛序列映到强收敛序列. 令 $\{h_{n}\}$ 为有界序列. 定理\cref{theorem:ch16-weak-seq}告诉我们 $\{h_{n}\}$ 有弱收敛的子序列. 该弱收敛子序列象强收敛.
\end{proof}

\begin{theorem}[Schauder]\label{theorem:ch16-schauder}
Hilbert 空间上的紧线性算子有紧的伴随.
\end{theorem}

\begin{proof}
令 $K\in\mathcal{L}(H)$ 是紧的. 根据前一个命题, 仅须证明 $K^{*}$ 将弱收敛序列映到强收敛序列. 令在 $H$ 中 $\{h_{n}\}\rightharpoonup h$. 对每个 $n$,
\begin{align}
\|K^{*}(h_{n})-K^{*}(h)\|^{2}=\langle KK^{*}(h_{n})-KK^{*}(h),h_{n}-h\rangle.\label{align:ch16-schauder-1}
\end{align}
由于 $K^{*}$ 是连续的, $\{K^{*}(h_{n})\}$ 弱收敛于 $K^{*}(h)$. 前一个命题告诉我们在 $H$ 中强意义下 $\{KK^{*}(h_{n})\}\to KK^{*}(h)$. 根据命题\cref{proposition:ch16-weak-properties},
\begin{align*}
\lim_{k\to\infty}\langle KK^{*}(h_{n})-KK^{*}(h),h_{n}-h\rangle=0.
\end{align*}
我们从\eqref{align:ch16-schauder-1}推出 $\{K^{*}(h_{n})\}$ 强收敛于 $K^{*}(h)$.
\end{proof}

\begin{exercise}\label{exercise:ch16-44}
证明: 若 $H$ 是无穷维的而 $T\in\mathcal{L}(H)$ 是可逆的, 则 $T$ 不是紧的.
\end{exercise}

\begin{exercise}\label{exercise:ch16-45}
证明命题\cref{proposition:ch16-compact-equiv}.
\end{exercise}

\begin{exercise}\label{exercise:ch16-46}
令 $\mathcal{K}(H)$ 表示 $\mathcal{L}(H)$ 中的紧算子子集. 证明 $\mathcal{K}(H)$ 是 $\mathcal{L}(H)$ 的线性子空间. 此外, 证明对 $K\in\mathcal{K}(H)$ 与 $T\in\mathcal{L}(H)$, $K\circ T$ 与 $T\circ K$ 都属于 $\mathcal{K}(H)$.
\end{exercise}

\begin{exercise}\label{exercise:ch16-47}
证明线性算子 $T:H\to H$ 是连续的当且仅当它将弱收敛序列映到弱收敛序列.
\end{exercise}

\begin{exercise}\label{exercise:ch16-48}
证明 $K\in\mathcal{L}(H)$ 是紧的当且仅当只要在 $H$ 中, $\{u_{n}\}\rightharpoonup u$ 且 $\{v_{n}\}\rightharpoonup v$, 就有 $\langle K(u_{n}),\ v_{n}\rangle\to\langle K(u),\ v\rangle$.
\end{exercise}

\begin{exercise}\label{exercise:ch16-49}
令 $\{P_{n}\}$ 为 $\mathcal{L}(H)$ 中的正交投影序列, 具有性质: 对自然数 $n$ 与 $m$, $P_{n}(H)$ 和 $P_{m}(H)$ 是 $H$ 的正交有限维子空间. 令 $\{\lambda_{n}\}$ 为实数的有界序列. 证明:
\begin{align*}
K=\sum_{n=1}^{\infty}\lambda_{n}\cdot P_{n}
\end{align*}
是恰当定义的 $\mathcal{L}(H)$ 中的对称算子, 它是紧的当且仅当 $\{\lambda_{n}\}$ 收敛到 $0$.
\end{exercise}

\begin{exercise}\label{exercise:ch16-50}
$X$ 是 Banach 空间, 定义算子 $T\in\mathcal{L}(X)$ 为紧的, 若 $T(B)$ 有紧闭包. 证明对一般的 Banach 空间命题\cref{proposition:ch16-compact-equiv}成立, 而对自反的 Banach 空间命题\cref{proposition:ch16-compact-weak-strong}成立.
\end{exercise}

\end{document}